% P16: finite longitudinal interventions, g-computation and policy weighting.
\section{Sequential interventions and the longitudinal g-formula}
\label{sec:p16-sequential-g-formula}

Hern\'an and Robins distinguish treatment strategies, several strengths of
sequential exchangeability, and the three representations of the g-formula
\cite[Chs.~19--20, Secs.~21.1--21.2 and Fine Point~21.2,
pp.~255--285, 291]{HernanRobins2020WhatIf2026}.
The version used here is dated 19 August 2026. This finite treatment supplies
the policy random variables and complete identification proofs. The source's
outcome-only argument and its joint trajectory argument have different
assumptions and are both retained. All histories below are chronological;
the source's more general topological ordering is not silently substituted
for the information available when a decision is made.

\subsection{A longitudinal unit and a policy}

\begin{hbdefinition}{def:p16-longitudinal-table}{Finite histories and consistent potential trajectories}
Fix an integer \(K\geq0\), finite nonempty sets \(\mathcal L_k\),
\(\mathcal A_k\) for \(0\leq k\leq K\), and a finite nonempty real
response set \(\mathcal Y\). The observed variables, in order, are
\[
 O=(L_0,A_0,L_1,A_1,\ldots,L_K,A_K,Y).
\]
Put \(\bar L_k=(L_0,\ldots,L_k)\),
\(\bar A_k=(A_0,\ldots,A_k)\), and
\(H_k=(\bar L_k,\bar A_{k-1})\). A history ending at \(-1\)
is an empty tuple; empty products equal one. Time indices in products
range from zero unless otherwise specified. Identify \(H_0\) with \(L_0\)
by dropping its empty action tuple. A history value is
\(h_k=(\bar l_k,\bar a_{k-1})\);
\((h_k,a_k,l_{k+1})\) denotes the corresponding next history.

Let \((\Omega,q)\) be one finite probability space carrying the observed
variables with the stated value sets. On this same space specify, for
every fixed treatment sequence
\(\bar a=(a_0,\ldots,a_K)\), potential variables
\[
 L_0,\quad L_k[\bar a_{k-1}]\ (1\leq k\leq K),\qquad Y[\bar a].
\]
Indexing the covariates only by earlier actions expresses no anticipation.
Require sequential consistency pointwise, with \(1\leq k\leq K\)
in the covariate identity:
\[
 \bar A_{k-1}=\bar a_{k-1}\Longrightarrow
 L_k[\bar a_{k-1}]=L_k,
 \qquad \bar A_K=\bar a\Longrightarrow Y[\bar a]=Y.
\]
The potential variables for all fixed sequences share this one outcome
space; no independence between worlds is required. Their complete finite
tuple is denoted \(\mathcal T\). There are
\(\prod_{k=0}^K|\mathcal A_k|\) static sequences, in particular
\(2^{K+1}\) when every treatment is binary.
\index{longitudinal history}\index{sequential consistency}
\end{hbdefinition}

\begin{hbdefinition}{def:p16-longitudinal-policy}{Static, dynamic and stochastic policies}
A known policy \(\pi=(\pi_0,\ldots,\pi_K)\) gives a normalized
mass function \(\pi_k(\cdot\mid h_k)\) on \(\mathcal A_k\) at
every formal history. It uses only that history. A deterministic dynamic
policy \(g\) has \(\pi_k(a\mid h)=\mathbf1_{\{a=g_k(h)\}}\).
A static policy fixes \(g_k(h)=a_k^*\); a stochastic policy permits
nondegenerate masses, with static stochastic policies independent of the
covariate history.

To define its counterfactuals, enlarge the space by independent finite
variables \(R_{k,h}\), one per time and formal history, with mass
\(\pi_k(\cdot\mid h)\). They are mutually independent and independent
of every original variable, including \(\mathcal T\). Define
\(H_k^\pi=(\bar L_k^\pi,\bar A_{k-1}^\pi)\) and evaluate recursively
\[
\begin{aligned}
 L_0^\pi&=L_0,& A_k^\pi&=R_{k,H_k^\pi},\\
 L_{k+1}^\pi&=L_{k+1}[\bar A_k^\pi]\quad(k<K),&
 Y^\pi&=Y[\bar A_K^\pi].
\end{aligned}
\]
Here a random index means evaluation of the already specified finite table
at its attained index. Write \(O^\pi\) for the resulting complete trajectory.
For deterministic \(g\), no extra randomness is needed and use the superscript
\(g\). The policy value is \(\psi(\pi)=\mathbb E[Y^\pi]\);
an effect contrast specifies two policies, for example
\(\psi(\pi)-\psi(\pi')\). A rule depending on the current natural treatment
before replacement, an unobserved covariate, or a future variable is not a
policy of this definition.
\index{treatment strategy!dynamic}\index{stochastic intervention}
\end{hbdefinition}

\begin{hbtheorem}{thm:p16-sequential-policy-semantics}{proved=annex:p16.sequential-policy-semantics}{Policy solutions, consistency and elimination of past action arguments}
The recursion defines a unique trajectory at every enlarged outcome.
On the event that factual actions agree with its selected policy actions
at every factual history, \(O^\pi=O\). Agreement through time \(k-1\)
already implies agreement of histories through \(H_k\).

For deterministic \(g\), recursively substitute earlier prescribed actions
into later rules to obtain functions \(g'_k(\bar l_k)\) with no action
arguments. The two policies give identical trajectories pointwise, and a
factual history follows one through a given time exactly when it follows
the other. In particular deterministic rules depending only on past
actions generate a single static sequence from the start of follow-up.
\end{hbtheorem}

\subsection{Exchangeability is an assumption about counterfactuals}

\begin{hbdefinition}{def:p16-sequential-exchangeability}{Outcome-specific and joint sequential exchangeability}
For deterministic \(g\), a history \(h_k\) is \emph{compatible} if
\(a_j=g_j(h_j)\) for all \(j<k\). Outcome-specific exchangeability for
this policy means, for every positive compatible observed history,
\[
 Y^g\mathrel{\perp\!\!\!\perp} A_k\mid H_k=h_k.
 \tag{SE--Y}
\]
Conditional independence at a history denotes factorization of the finite
conditional joint law there. The condition is imposed only where the
history has positive probability.

For fixed \(\bar a\), put
\[
 T_k(\bar a)=
 (L_{k+1}[\bar a_k],\ldots,L_K[\bar a_{K-1}],Y[\bar a]);
\]
for \(k=K\) this tuple consists just of \(Y[\bar a]\).
Joint sequential exchangeability means, for every \(\bar a,k\) and every
positive \(h_k=(\bar l_k,\bar a_{k-1})\),
\[
 T_k(\bar a)\mathrel{\perp\!\!\!\perp} A_k\mid H_k=h_k.
 \tag{SE--J}
\]
This is joint independence of one fixed world's future covariates and
response. It does not assert independence of the joint collection of all
worlds from treatment. Requiring (SE--Y) only for static policies is weaker
than requiring it for all dynamic policies. These distinctions follow the
source's Technical Points~19.2--19.3, pp.~261--262.
\end{hbdefinition}

\begin{hbtheorem}{thm:p16-sequential-exchangeability-bridge}{proved=annex:p16.sequential-exchangeability-bridge}{From fixed-world joint exchangeability to dynamic policies}
Under sequential consistency, (SE--J) implies, for every deterministic
\(g\) and positive compatible \(h_k\),
\[
 (L_{k+1}^g,\ldots,L_K^g,Y^g)
 \mathrel{\perp\!\!\!\perp} A_k\mid H_k=h_k.
\]
Conversely, these joint dynamic-policy conditions include (SE--J) by
choosing static policies. Thus the two families of joint conditions are
equivalent. Independence of the response and future covariates separately
cannot replace their joint independence in this assertion.
\end{hbtheorem}

\begin{hbtheorem}{thm:p16-sequential-structural-bridge}{proved=annex:p16.sequential-structural-bridge}{A sufficient structural assignment mechanism}
Consider the existing finite structural model in the order of \(O\), with
additional background \(U\). Suppose every covariate and response equation
uses its earlier observed variables and \(U\), and treatment equations are
\[
 A_k=f_k(H_k,E_k),
\]
where \(E_0,\ldots,E_K,U\) are mutually independent. Each \(E_k\)
is used directly only by its own treatment equation. Components of \(U\)
may be dependent and shared by covariate and response equations.
Then fixed equation replacement supplies the potential table and sequential
consistency, and
\(\mathcal T\mathrel{\perp\!\!\!\perp} A_k\mid H_k\).
Consequently (SE--J) holds. Replacing treatment equations by the policy
selectors gives exactly \(O^\pi\).

These conditions describe a sequential randomization mechanism based on
the recorded history. An arbitrary structural model, or knowledge of an
observed treatment conditional distribution alone, does not imply them.
\end{hbtheorem}

\subsection{The exact support domain}

\begin{hbdefinition}{def:p16-sequential-observed-kernels}{Observed kernels and reachable-history positivity}
Let \(p\) be the full law of \(O\). On positive conditioning events define
\[
\begin{aligned}
 b(l_0)&=p(L_0=l_0),\\
 e_k(a\mid h)&=p(A_k=a\mid H_k=h),\\
 r_k(l\mid h,a)&=p(L_{k+1}=l\mid H_k=h,A_k=a)\quad(k<K),\\
 r_K(y\mid h,a)&=p(Y=y\mid H_K=h,A_K=a).
\end{aligned}
\]
At other conditioning arguments choose any normalized versions of these
kernels. Such a choice specifies numbers, not a conditional probability
at a null event.

Construct masses \(m_0(l_0)=b(l_0)\). At each \(m_k(h)>0\), permit
only actions with \(\pi_k(a\mid h)>0\), require
\[
 p(H_k=h)>0,\qquad e_k(a\mid h)>0,                 \tag{RP}
\]
and propagate, for \(k<K\),
\[
 m_{k+1}(h,a,l)=m_k(h)\pi_k(a\mid h)r_k(l\mid h,a).
\]
Failure of either required positive denominator is a failed support
condition, before evaluating a child kernel there. Histories with zero
propagated mass and actions with zero policy probability create no branch.
This is \emph{reachable-history positivity for the g-computation tree}.
The equivalent observed-law condition is: at every positive observed
history with \(\prod_{j<k}\pi_j(a_j\mid h_j)>0\), every action of positive
\(\pi_k\)-probability has positive \(e_k\)-probability.
Under the joint identification theorem, \(m_k\) is the actual policy
history law, so (RP) is precisely positivity on policy-reachable histories.
\index{positivity!reachable history}
\end{hbdefinition}

\begin{hblemma}{thm:p16-sequential-support}{proved=annex:p16.sequential-support}{Propagation, normalization and independence of null versions}
Under (RP), every reached history and action has a positive observed
conditioning event. The propagated masses sum to one at each time, and
\[
 q_\pi(o)=b(l_0)\left\{\prod_{k=0}^K\pi_k(a_k\mid h_k)\right\}
             \left\{\prod_{k=0}^{K-1}r_k(l_{k+1}\mid h_k,a_k)\right\}
             r_K(y\mid h_K,a_K)                         \tag{GF}
\]
defines a normalized law independent of all null-event versions. A product
on a branch already assigned zero mass is evaluated as zero, without
querying any later conditional. Also \(q_\pi(o)>0\) implies \(p(o)>0\).
The displayed observed-law positivity condition is equivalent to (RP).
These are probability-algebra conclusions and do not yet assert causality.
\end{hblemma}

\subsection{Two causal justifications of g-computation}

\begin{hbtheorem}{thm:p16-sequential-joint-g-formula}{proved=annex:p16.sequential-joint-g-formula}{The full policy trajectory law}
Under sequential consistency, (SE--J), and (RP) for the specified known
policy, \(\mathbb P(O^\pi=o)=q_\pi(o)\). Therefore
\[
 \mathbb P(Y^\pi=y)=\sum_{\bar l,\bar a}q_\pi(\bar l,\bar a,y),
 \qquad \psi(\pi)=\sum_{\bar l,\bar a,y}yq_\pi(\bar l,\bar a,y).
\]
All sums are finite. For deterministic \(g\), only the recursively
prescribed actions \(a_k=g_k(\bar l_k,\bar a_{k-1})\) contribute; for
static \(\bar a^*\), the formula becomes
\[
\begin{aligned}
 \mathbb E[Y[\bar a^*]]
 &=\sum_{\bar l}\mathbb E[Y\mid H_K=(\bar l,\bar a^*_{K-1}),A_K=a_K^*]\\
 &\hspace{3em}\cdot b(l_0)\prod_{k<K}r_k(l_{k+1}\mid h_k,a_k^*),
\end{aligned}
\]
where \(h_k=(\bar l_k,\bar a^*_{k-1})\), with the same zero-branch
convention. The policy law is identified among
all finite potential-table models sharing \(p\) and satisfying these
assumptions, for the stated fixed value sets and policy.

Under consistency and (SE--J), the operational condition (RP) is equivalent
to requiring \(p(H_k=h)>0\) and \(e_k(a\mid h)>0\) whenever
\(\mathbb P(H_k^\pi=h)\pi_k(a\mid h)>0\). Thus no action outside the
actual reachable policy support is needed for this theorem.
\end{hbtheorem}

\begin{hbtheorem}{thm:p16-sequential-outcome-g-formula}{proved=annex:p16.sequential-outcome-g-formula}{The distinct outcome-only identification route}
For a specified deterministic \(g\), sequential consistency, (SE--Y)
for that \(g\), and (RP) suffice for
\[
 \mathbb E[\varphi(Y^g)]=\sum_o\varphi(y)q_g(o)
\]
for every real function \(\varphi\) on \(\mathcal Y\). In particular the
outcome law and mean are identified. The conclusion concerns the response;
it supplies no equality between the policy covariate-trajectory law and
the covariate marginal of \(q_g\). Indeed the premises can hold in two
models with identical observed laws and different intervened covariate
laws. It applies to a static policy without assuming exchangeability for
other policies.
\end{hbtheorem}

\begin{hbtheorem}{thm:p16-sequential-ice}{proved=annex:p16.sequential-ice}{Iterated conditional expectations}
For a specified \(\varphi\), define on reached histories and selected actions
\[
\begin{aligned}
 B_K(h,a)&=\sum_y\varphi(y)r_K(y\mid h,a),\\
 V_K(h)&=\sum_a\pi_K(a\mid h)B_K(h,a),\\
 B_k(h,a)&=\sum_l r_k(l\mid h,a)V_{k+1}(h,a,l),\\
 V_k(h)&=\sum_a\pi_k(a\mid h)B_k(h,a)\quad(k<K).
\end{aligned}
\]
Terms of zero transition or policy mass are omitted before evaluating their
continuation. Then
\[
 \sum_{l_0}b(l_0)V_0(l_0)=\sum_o\varphi(y)q_\pi(o).
\]
This recursion terminates after \(K+1\) backward stages. Its right side
equals \(\mathbb E[\varphi(Y^\pi)]\) under the joint theorem, and for
deterministic \(g\) also under the outcome-only theorem. For \(k<K\),
\(B_k(h,a)\) is the observed conditional expectation of the specified
function \(V_{k+1}(H_{k+1})\) given \(H_k=h,A_k=a\). It is not obtained
by changing the observed law inside an unqualified conditioning expression.
\end{hbtheorem}

\begin{hbtheorem}{thm:p16-sequential-ipw}{proved=annex:p16.sequential-ipw}{The known-policy inverse-probability identity}
On an observed assignment with \(p(o)>0\), put
\[
 W^\pi(o)=\prod_{k=0}^K\frac{\pi_k(a_k\mid h_k)}{e_k(a_k\mid h_k)};
\]
put \(W^\pi(o)=0\) on assignments with \(p(o)=0\). Under (RP),
\[
 p(o)W^\pi(o)=q_\pi(o),\qquad
 \mathbb E_p[W^\pi]=1,\qquad
 \mathbb E_p[W^\pi\varphi(Y)]=\sum_o\varphi(y)q_\pi(o).
\]
Under the corresponding causal hypotheses this last quantity is
\(\mathbb E[\varphi(Y^\pi)]\). For deterministic \(g\), its numerator
is the indicator that every factual action follows \(g\).

For static \(\bar a^*\), multiplying this indicator weight by a known
positive constant \(c(\bar a^*)\) leaves the normalized population ratio
unchanged. In particular a product of treatment-only stabilization
probabilities is such a constant on a fixed static sequence. Under a
dynamic rule its value can vary between compatible histories, and the
same cancellation is generally invalid.

If \(\pi_k=e_k\) at every positive observed history, (RP) holds and
\(q_\pi=p\). With the causal hypotheses this says that independent
policy reassignment using the same history-specific probabilities preserves
the population trajectory law. It need not preserve an individual unit's
realized trajectory.
\end{hbtheorem}

\subsection{Calculated treatment-confounder feedback}

\begin{hbtheorem}{thm:p16-sequential-source-feedback}{proved=annex:p16.sequential-source-feedback}{The source's null example with a complete finite realization}
Treat the source's Table~20.1, repeated as Table~21.1, as a stipulated
population table, with each count divided by \(32000\). Its entries are
\[
\begin{array}{ccc|r|r}
 A_0&L_1&A_1&32000p(A_0,L_1,A_1)&\mathbb E[Y\mid A_0,L_1,A_1]\\\hline
 0&0&0&2400&84\\0&0&1&1600&84\\
 0&1&0&2400&52\\0&1&1&9600&52\\
 1&0&0&4800&76\\1&0&1&3200&76\\
 1&1&0&1600&44\\1&1&1&6400&44
\end{array}
\]
There is a finite sequentially randomized structural realization with
\(Y\in\{20,100\}\) unaffected by either treatment. It has
\(\mathbb E[Y^\pi]=60\) under every policy.
G-computation for never treating gives \(84/4+3(52)/4=60\);
for always treating it gives \((76+44)/2=60\). IP weighting gives
the same two results. The unadjusted always-versus-never contrast is
\(164/3-68=-40/3\); the contrast within either level of \(L_1\)
is \(-8\). Neither is the null policy contrast.
The annex specifies the entire outcome distribution, every structural
randomization probability, and the source table's weighting calculation.
\end{hbtheorem}

\begin{hbtheorem}{thm:p16-sequential-feedback-effects}{proved=annex:p16.sequential-feedback-effects}{Feedback with nonzero static, dynamic and stochastic effects}
Let \(L_0\) be constant, \(A_0\) fair, and let \(E_L,E_A\) be
independent Bernoulli\((1/4)\) bits independent of \(A_0\). Define
\[
 L_1=A_0\mathbin\oplus E_L,\qquad
 A_1=L_1\mathbin\oplus E_A,\qquad Y=L_1+2A_1.
\]
The complete observed law has eight nonzero cells:
\[
\begin{array}{ccc|c|c}
 A_0&L_1&A_1&Y&32p\\\hline
 0&0&0&0&9\\0&0&1&2&3\\0&1&0&1&1\\0&1&1&3&3\\
 1&0&0&0&3\\1&0&1&2&1\\1&1&0&1&3\\1&1&1&3&9
\end{array}
\]
All four static means are \(1/4+a_0/2+2a_1\). Fixing \(A_0=a\)
and then treating exactly when \(L_1=1\) gives mean
\(3/4+3a/2\). Replacing that second rule by
\(\pi_1(1\mid h)=1/5+3L_1/5\) gives mean
\(19/20+11a/10\). G-computation, the backward recursion, and known-policy
weighting give these same values. The always-versus-never causal contrast
is \(5/2\), the unadjusted contrast is \(14/5\), and comparing within a
fixed \(L_1\) gives \(2\).
\end{hbtheorem}

\subsection{Failures of assumptions and changes of target}

\begin{hbtheorem}{thm:p16-sequential-failures}{proved=annex:p16.sequential-failures}{Positivity and exchangeability cannot simply be removed}
The following finite counterexamples satisfy sequential consistency.
\begin{enumerate}
\item With two treatment times and constant covariates, let \(A_0\) be fair
and \(A_1=0\). The models \(Y=0\) and \(Y=A_1\) have identical
observed laws and satisfy (SE--J), but \(\mathbb E[Y^{(0,1)}]\) is
respectively zero and one. The requested second action fails (RP).
The policy \((0,0)\) is supported and has identified mean zero in both.
\item With one treatment time and constant covariates, let \(U\) be fair,
\(E\) an independent Bernoulli\((1/4)\) bit, and \(A=U\). The models
\(Y=A\mathbin\oplus U\mathbin\oplus E\) and \(Y=E\) share the strictly
positive observed law \(p(a,0)=3/8,p(a,1)=1/8\), but under either fixed
action their response probabilities are respectively \(1/2\) and \(1/4\).
The first model violates exchangeability. Thus treatment positivity alone
does not identify the policy value.
\item Static outcome-specific exchangeability, even at all static sequences,
does not imply dynamic outcome-specific exchangeability. There is a
two-time example satisfying every observed treatment positivity condition,
with separately exchangeable fixed-world future covariate and response,
but a dynamic response that is not exchangeable. Its construction and
unequal g-formula and true dynamic means are given in the annex.
\end{enumerate}
These examples establish failure in their specified classes. Failure of
one sufficient exchangeability condition does not assert universal
nonidentification by every other available causal argument.
\end{hbtheorem}

\begin{hbtheorem}{thm:p16-sequential-weight-boundaries}{proved=annex:p16.sequential-weight-boundaries}{Dynamic stabilization and fitted denominators change the calculation}
Let \(K=0\), \(L\) be fair, \(A\) independent with success probability
\(3/4\), and \(Y=L\). Under the dynamic rule \(A=L\), the policy value
is \(1/2\). The normalized mean among adherers weighted by
\(p(A)/p(A\mid L)\) is instead \(3/4\). The correct indicator divided
by \(p(A\mid L)\) gives \(1/2\).

If \(O_1,\ldots,O_n\) are identically distributed observations from the
stated population, the average of \(W^\pi(O_i)Y_i\) using the true
\(e_k\) has expectation \(\psi(\pi)\) under the causal hypotheses.
Independence between observations is unnecessary for that expectation
identity. Replacing \(e_k\) by fitted probabilities does not preserve
the identity automatically; even fixed incorrect denominators can fail.
No sample-size limit, variance formula, model-fitting guarantee or empirical
treatment recommendation follows from the population identification theorem.
\end{hbtheorem}

The three population representations concern a known policy and the actual
observed probability law. Estimation by predictive models, doubly robust
procedures, censoring interventions, continuous conditional kernels,
optimal-policy learning and modified treatment policies require their own
assumptions and results. General ID and IDC address different graphical
identification questions; their success is not a premise in the proofs here.
